\chapter*{Acknowledgments}
\addcontentsline{toc}{chapter}{Acknowledgments}

First and foremost, I would like to thank my supervisors,
Prof. António Paulo Moreira and Héber Miguel Sobreira,
for their unwavering support throughout my PhD journey and for the opportunity
to work with them. This journey would not have been possible without them,
including the invitation by Prof.~António Paulo upon completing my M.Sc.
dissertation, starting to work with Héber at \CRIIS from \INESCTEC,
and then exploring and researching perception,
Simultaneous Localization and Mapping~(SLAM), calibration, control systems, and
mobile robotics in general.

Next, a special thanks to my colleagues from CRIIS,
including from the laboratories \iilab and \TRIBElab.
Although my research is related to the intralogistics industry and more
related to iilab's scope and mission,
I began my PhD journey working side by side with TRIBElab, as the lab is located
on the campus of the Faculty of Engineering, University of Porto~(FEUP).
This experience was excellent in terms of the team environment and spirit,
while also providing the opportunity to interact daily with my mates from the
master's degree, namely André and Daniel. Additionally, the TRIBElab leader,
Filipe Santos, was one of the most important people who helped me decide that
staying at INESC~TEC would be the best choice to continue my journey and pursue
a PhD degree, given the interaction with real-world applications and the
opportunity to move beyond simulators and the theoretical world.
So, a special thanks to him for that support, and to André and Daniel,
two close companions and friends who supported me greatly throughout my PhD.
After I moved to the iilab facilities, it has been a fantastic journey,
marked by the creation of new team partnerships and demonstrating that working
side by side with colleagues from the same application field and different
research areas can result in fascinating research experiences.
A special thanks to my team leader Héber, also the co-supervisor of my
PhD thesis, Prof. Manuel Silva, for his guidance on the GreenAuto project
directly linked with my PhD research and all the support given by him and the
current and past CRIIS coordination team, to all my teammates of the
Mobile Robotics Development Team~(MRDT), and all my colleagues from iilab that
I interact daily with and has been a pleasure to work with them.

Another very special person in the PhD program was Giorgio~Grisetti from
Sapienza University of Rome. Upon gathering persons of interest for a PhD
internship, Giorgio clearly distinguished himself by having everything I wanted
to supervise me in that experience: theoretical background, both in filter and
graph-based formulations for the SLAM problem, and joining the theory with a
tremendous experience in implementing SLAM algorithms with real-world
applications, including having developed countless open-source applications used
by everyone in the world (GMapping, collaboration in the g2o framework, among
others). I contacted him by email on July~18th, 2023, a Tuesday, after dinner,
thinking that I would only receive an answer after the summer vacations or not
at all. To my complete surprise, the very next day, at 8~AM (Rome local time),
I received a meeting invitation from Giorgio scheduled for Friday of the same
week. At that meeting, the internship was agreed upon, including the timeframe
(common availability between January and June 2024 and a time period of three
months), the scope, and, especially, how to prepare for the internship, by doing
the entire Probabilistic~Robotics course from Giorgio that had online video
recordings from the COVID time (2020/21 and 2021/22 academic years).
Indeed, the experience in Rome was life-changing in terms of my research,
having a tremendous impact on my PhD thesis, as evidenced by the extension of my
stay from three to six months (from February until July 2023). Of course, a
lot of that experience was thanks to Giorgio, for being the great person he is,
not caring about societal status (in terms of academic background and
professional success) and being hands-on, interacting daily with his PhD
students and with external ones like me. There are no words enough to thank him
for that, a massive \textbf{Thank You} for the experience!
Additionally, I would like to thank my colleagues from the
Robotics Vision and Perception Group (RVP-Group) and the LabRoCoCo for their
hospitality and engaging research discussions. Last but not least, a big thank
you to INESC~TEC and, especially, to the past CRIIS coordination team
(Prof.~António~Paulo and Germano~Veiga) for giving me the opportunity to
undertake the PhD internship. Although I transitioned from the national
Fundação para a Ciência e a Tecnologia~(FCT) PhD scholarship to a researcher
contract at INESC~TEC in April~2023, the CRIIS coordination not only supported
but also encouraged me to go to the internship. So, a big thank you to them and
the teammates at CRIIS, who, of course, supported me throughout the journey.

In parallel with my PhD, I had the opportunity to be part of the 5dpo Robotics
Team at FEUP. I want to thank all past and present members of the team, with
special thanks to Prof.~António~Paulo, Cláudia~Rocha,
João~Martins, José~Pedro, José~Sarmento,
and Maria~Lopes. Thank you for all the time we spent together, for our
participation in the Portuguese Robotics Open~(FNR), for the frustration of the
defeats that helped us grow as a team, and for the happiness, experience,
knowledge, and pride that came with our victories. Thank you also for all the
hard work, partnership, friendship, and the many engaging discussions, forming
the backbone of an incredible team and experience, one in which I truly felt
part of a team fighting toward common goals, working together without hidden
interests, supporting one another to achieve our objectives while enjoying our
time together and learning from each other.

Fortunately, I have always had a strong set of friendships behind me, supporting
me since high school, my master's degree, and now throughout my PhD. A big thank
you to my friends from high school and further: Bárbara and Ruben, Bruno and
Juliana, Beatriz and Rocha, Hugo and Márcia, Mariana and Tommy, Pimentas, and
Sâncio. They are to me like a family, always having the weekly coffee in Gaia
to discuss economics, politics, a little bit of football (unfortunately, most of
them support Porto, not my Benfica!) and gossip. Furthermore, another big thanks
to all my friends from the masters at FEUP who are still present for the monthly
craft beer at Cervejaria do Carmo or Cervejaria Gulden Draak Bierhuis in Porto
and possibly dinner afterwards, depending on the mood, and the yearly Christmas
dinner, including Bruno, Daniel and Clara, Hélder, Mariano, Mariana and Granja,
Pedro, Ricardo and Isabel, Tiago, Tovar, and Xavi. It's quite fascinating that,
even after completing our master's degrees at FEUP, we have still been meeting
regularly since 2020. I will never forget the special moment when Bruno, Hélder,
Mariano, and Tovar took a road trip from Germany to Rome, Italy, and back,
during May 2024, while I was still on my PhD internship. Indeed, that weekend we
spent together exploring Rome and enjoying good food and drink was the best of
my whole stay in Italy, forever etched in my memory. Also, I would like to
extend another special thanks to Daniel and Pedro. We started as classmates,
study buddies, and now I consider you both as brothers. I thank you for all the
advice, discussions, and sharing of experiences that we have been doing, and I
hope to continue doing it throughout many years with you!

A special thank you to all the people who were important and had an impact on me
throughout my PhD journey in ways not previously mentioned. First and foremost,
to Daniela, who was and will continue to be, regardless of where life leads,
a truly special person to me. She supported me profoundly on a personal level,
and I am grateful for all the memories we shared, the advice we exchanged,
and the many long and enjoyable walks we took over these years.
Additionally, I would like to thank Helena and Joana for the valuable advice
they provided at various points along the way, demonstrating that both
physical and mental strength are essential to achieving success in life.
Moreover, I would also like to extend my gratitude to professors António~Paulo
and João~Tasso~Sousa, Ana~Rufino~Ferreira, Catarina~Pinto, and Héber for all the
support they provided during the preparation for applying to the dual degree PhD
program in Electrical and Computer Engineering from
Carnegie Mellon University~(CMU) at the end of 2020.
Although the application was not successful, their help was essential in that
endeavour, from providing recommendation letters to CV preparation,
drafting motivation letters, and searching for potential supervisors
on the CMU side. Their encouragement motivated me to pursue a PhD internship
abroad, which led me to Rome, where I had the great privilege of
working with Giorgio.

Last but not least, my deepest and most sincere thanks go to my family, my
parents, José Fernando and Maria, my brother Rafael and his partner.
Their support has been constant throughout my life. They have always been there
for me, ready to help in any way I needed, and they gave me the freedom to
explore, take risks, and make mistakes that helped me grow.
Work may come and go, as do friendships and professional connections, but
what remains unwavering is their presence. For that, I am endlessly grateful.

\vspace{10mm}
\flushleft{Ricardo B. Sousa}

\cleardoublepage
